
\documentclass[a4paper,12pt,oneside]{article} %einseitiger Druck
% Optionen:
% - a4paper => DIN A4-Format
% - 12pt    => Schriftgröße (weitere  
%         grundlegende Fontgrößen: 10pt, 11pt)
% - oneside => Einseitiger Druck

%% Verwendete Pakete:
\usepackage[ngerman]{babel} % für die deutsche Sprache
\usepackage{caption} % Für die Titel in Gleitobjekten (wie figure)
\usepackage[T1]{fontenc} % Für Sonderzeichen u.a.
\usepackage[utf8]{inputenc} % Für die direkte Eingabe von Umlauten
\usepackage{fancyhdr} % Für Kopf- und Fußzeilen
\usepackage{lscape} % Für Querformat
\usepackage{lmodern} % Type1-Schriftart für nicht-englische Texte
\usepackage{ulem} % Für das Unterstreichen von Text z.B. mit \uline{}
\usepackage[left=3cm,right=2cm,top=1.5cm,bottom=1cm,textheight=245mm,textwidth=160mm,includeheadfoot,headsep=1cm,
			footskip=1cm]{geometry} % Einrichtung der Seite 

\usepackage{graphicx} % Zum Laden von Graphiken

% Pakete für Tabellen
\usepackage{tabularx} % Einfache Tabellen
\usepackage{longtable} % Tabellen als Gleitobjekte (für die Aufteilung bei langen 
 %Tabellen über mehrere Seiten)
\usepackage{multirow} % Für das Verbinden von Zeilen innerhalb einer Tabelle mit
 % \multirow{anzahl}{*}{Text}

% Pakete für Formeln
\usepackage{amsmath}
\usepackage{amsthm}
\usepackage{listings}
\usepackage{amsfonts}
\usepackage{setspace} % Paket zum Setzen des Zeilenabstandes
\usepackage[colorlinks,pdfpagelabels,pdfstartview=FitH,
bookmarksopen=true,bookmarksnumbered=true,linkcolor=black,
plainpages=false,hypertexnames=false,citecolor=black]{hyperref} 
% Globale Einstellungen (gelten für das gesamte Dokument, sofern sie nicht verändert werden)
\setlength{\headheight}{14.599pt} % Kopzeilenhöhe festlegen
\setlength{\headwidth}{160mm}     % Kopzeilenbreite festlegen
\onehalfspacing % Zeilenabstand auf 1,5-zeilig setzen

 
%%%%%%%%%%%%%%%%%%%%%%%%%%%%%%%%%%%%%%%%%%%%%
%% DOKUMENT                                %%
%%%%%%%%%%%%%%%%%%%%%%%%%%%%%%%%%%%%%%%%%%%%%
\begin{document}
 
  % TO-DO: Unbeschriftetes Vorblatt (Leere Seite)
  % ...
   
  % Deckblatt 
  \pagestyle{empty}
  \section*{} %Sternchen damit die Seite nicht im Ihnhaltsverzeichnis auftaucht.
  
  \begin{titlepage}
    \includegraphics[scale=1.00]{Sources/logo_fhkoeln_sw}\\
    \begin{center}
      \Large
      Fachhochschule Köln\\
      Fakultät für Informatik und Ingenieurwissenschaften\\
      \hrule\par\rule{0pt}{2cm} % Horizontale Trennlinie  mit 2 cm Abtand nach unten erzeugen
      \LARGE
      \textsc{Documentation}\\
      \vspace{1cm} % Vertikaler Abstand von 1cm erzeugen
      \huge
      WPF Modern Web\\
      \Large
      \vspace{1.5cm}
      \large
      presented at Cologne University of Applied Science\\
      Campus Gummersbach\\
      Course of studies
      Medieninformatik\\ 
      \vspace{1.0cm}
      prepared by:\\
      \textsc{Dominik Schilling}\\
      (Matrikelnummer: ???)\\
      \textsc{Laura-Maria Hermann}\\
      (Matrikelnummer: 11083968)\\
      \textsc{Dario }\\
      (Matrikelnummer: ???)\\
      \vspace{1.5cm}
      \vspace{1.5cm}
     
    \end{center}    
  \end{titlepage}
  
  \newpage
 
  
  \tableofcontents
  % TO-DO: Seitennummerierung "einschalten"
  % ...
	\pagestyle{fancy}
   		
 \newpage %neue Seite zwischen Inhaltsverzeichnis und Abbildungsverzeichnis
  % TO-DO: Abbildungsverzeichnis einfügen
  % ...
  
 

  

\newpage

\section{Introduction}

This documentation contains the content elaborated within the scope 
of the WPF modern Web applications. It includes, among other things,
the source code and the approaches or considerations that have been 
made during the early stage of the project.
The goal was to create a gallery in form of a website, which had
to include the following features:
\begin{itemize}
\item The gallery as such should be implemented using Javascript. 
\item Individual images should be zoomed when selected and navigable with arrows. 
\item A log-in area should be available. 
\item MySQL implementation to allow log-in for users, but also to allow the upload of new images. 
\item JSON, JQuery and AJAX should be used.
\end{itemize}

\section{Planning}

The first task was to discuss how the gallery should look like
- not optically but functionally. Because a log-in area was required,
it seemed reasonable to create an administrator with extended
rights (although initially irrelevant). (?)

When looking at existing galleries on the net it is possible to 
roughly distinguish between two groups. 

On the one hand, there are web pages which serve as a real galleries and 
only have one author who can publish content. In this case only the admin 
- the author - himself can eg. upload, delete, edit, etc. photos.

On the other hand, many galleries are used as a platform.
These web pages give not only administrators, but also registered users 
the possibility to upload and manage their photos.

Due to the WPF we had to choose the second alternative. 


\section{Realization}


\subsection{MVC-paradigm}

The idea of administrators and users beeing able to upload and manage photos, 
was pursued throughout the entire implementation of the MVC paradigm.
In the following this will be further explained in more detail. 

\emph{What is MVC?}

MVC (Model-view-controller) is an software architecture pattern which provides a 
strict separation between reprasentation of information and user's interaction with it. 
A distinction is made between the three components model, view and controller. These
are explained in more detail below.


\subsubsection*{Request}

Example: 
\begin{lstlisting}
GET localhost/user/?foo=bar
\end{lstlisting}

The Request class extracts the URL into its component parts. User,
for example, is one segment.


\subsubsection*{Controller}

The controllers include the logic that always belongs to a resource. 

Example: 

\begin{lstlisting}
localhost / user
\end{lstlisting}
then the controller is loaded, see load.php. The controllers are
located in "includes / controllers" directory. If no other resource
is specified, for example "localhost / user / blaa"
the index method of the controller is called and will start. If the resource was
"blaa" and the indicated method exists, it is called. When it does not
exist, the "index" is retrieved.


\subsubsection*{Views}

The views are located in the 
includes / views
directory. The views are responsible for the visual representation
of the data.

Example call:

\begin{lstlisting}
$view = new View( 'install/success' ); 
$view->set_page_title( 'Installation' );  
$view->render(); // Gibt das Template aus
\end{lstlisting}

A view can consist of several templates, eg headers - content - footer.


\subsubsection*{Model}


\subsection{Database}

One of the first tasks was creating a database in which all uploaded
images are stored and can be retrieved. In the following, all the
contents and information related to the database are examined in detail.


\subsubsection{Create}

First of all, a database was created. Each user can later use the graphical
interface to create a gallery and thus create a set of tables in the
database. In the following explanations of all realized tables and fields
can be found. Collation: utf8\_general\_ci


\subsubsection{Database tables}


\paragraph*{Users}

This is to inform the users of the gallery. These are uniquely identified
by an ID.
\begin{tabular}{|c|c|}
\hline 
Feld & Typ\tabularnewline
\hline 
\hline 
ID & bigint(20)\tabularnewline
\hline 
user\_login & varchar(50)\tabularnewline
\hline 
user\_pass & varchar(64)\tabularnewline
\hline 
user\_email & varchar(100)\tabularnewline
\hline 
user\_registered & datetime\tabularnewline
\hline 
user\_status & int(10)\tabularnewline
\hline 
\end{tabular}


\paragraph*{Usermeta}

Metadatas of the users

\begin{tabular}{|c|c|}
\hline 
Feld & Typ\tabularnewline
\hline 
\hline 
ID & bigint(20)\tabularnewline
\hline 
user\_id & bigint(20)\tabularnewline
\hline 
meta\_key & varchar(255)\tabularnewline
\hline 
meta\_value & longtext\tabularnewline
\hline 
\end{tabular}

\paragraph*{Images}

Data of the images
\begin{tabular}{|c|c|}
\hline 
Feld & Typ\tabularnewline
\hline 
\hline 
ID & bigint(20)\tabularnewline
\hline 
user\_id & bigint(20)\tabularnewline
\hline 
uploaded\_date & datetime\tabularnewline
\hline 
image\_caption & longtext\tabularnewline
\hline 
image\_title & text\tabularnewline
\end{tabular}

\paragraph*{Imagemeta}

Metadatas of the images

\begin{tabular}{|c|c|}
\hline 
Feld & Typ\tabularnewline
\hline 
\hline 
ID & bigint(20)\tabularnewline
\hline 
image\_id & bigint(20)\tabularnewline
\hline 
metakey & varchar(255)\tabularnewline
\hline 
meta\_value & longtext\tabularnewline
\end{tabular}

\paragraph*{Gallery}

Informations about the gallery in general

\begin{tabular}{|c|c|}
\hline 
Feld & Typ\tabularnewline
\hline 
\hline 
ID & bigint(20)\tabularnewline
\hline 
user\_id & bigint(20)\tabularnewline
\hline 
public & boolean\tabularnewline
\hline 
gallery\_title & text\tabularnewline
\hline 
gallery\_description & longtext\tabularnewline
\end{tabular}

\paragraph*{ImageRelationships}

\begin{tabular}{|c|c|}
\hline 
Feld & Typ\tabularnewline
\hline 
\hline 
ID & bigint(20)\tabularnewline
\hline 
gallery\_id & bigint(20)\tabularnewline
\end{tabular}

\subsubsection{SQL-Query}

The following SQL commands are called to create the tables.

\begin{lstlisting}
CREATE TABLE `gallery` (
  `ID` bigint(20) unsigned NOT NULL,
  `user_id` bigint(20) unsigned NOT NULL,
  `is_public` tinyint(1) NOT NULL DEFAULT '1',
  `gallery_title` text NOT NULL,
  `gallery_description` longtext NOT NULL,
  PRIMARY KEY (`ID`)
) DEFAULT CHARSET=utf8;

CREATE TABLE `imagemeta` (
  `ID` bigint(20) unsigned NOT NULL AUTO_INCREMENT,
  `image_id` bigint(20) unsigned NOT NULL,
  `meta_key` varchar(255) NOT NULL,
  `meta_value` longtext NOT NULL,
  PRIMARY KEY (`ID`),
  KEY (`image_id`),
  KEY (`meta_key`)
) DEFAULT CHARSET=utf8;

CREATE TABLE `images` (
  `ID` bigint(20) unsigned NOT NULL AUTO_INCREMENT,
  `user_id` bigint(20) unsigned NOT NULL,
  `uploaded_date` datetime NOT NULL,
  `image_title` text NOT NULL,
  `image_description` longtext,
  `image_caption` text NOT NULL,
  PRIMARY KEY (`ID`),
  KEY (`user_id`)
) DEFAULT CHARSET=utf8;

CREATE TABLE `image_relationships` (
  `image_id` bigint(20) unsigned NOT NULL,
  `gallery_id` bigint(20) unsigned NOT NULL,
  PRIMARY KEY (`image_id`, `gallery_id`),
  KEY (`gallery_id`)
) DEFAULT CHARSET=utf8;

CREATE TABLE `usermeta` (
  `ID` bigint(20) unsigned NOT NULL AUTO_INCREMENT,
  `user_id` bigint(20) unsigned NOT NULL,
  `meta_key` varchar(255) NOT NULL,
  `meta_value` longtext NOT NULL,
  PRIMARY KEY (`ID`),
  KEY (`user_id`),
  KEY (`meta_key`)
) DEFAULT CHARSET=utf8;

CREATE TABLE `users` (
  `ID` bigint(20) unsigned NOT NULL AUTO_INCREMENT,
  `user_login` varchar(60) NOT NULL,
  `user_pass` varchar(64) NOT NULL,
  `user_email` varchar(100) NOT NULL,
  `user_registered` datetime NOT NULL DEFAULT '0000-00-00 00:00:00',
  `user_status` int(10) NOT NULL,
  PRIMARY KEY (`ID`),
  KEY (`user_login`),
  KEY (`user_email`)
) DEFAULT CHARSET=utf8;
\end{lstlisting}





\subsection{resources }

There follows an explanation of the necessary resources.


\subsubsection{/register}

At this point, the registration of new users was possible. This only
needs to be done if wanting to create a gallery and to post pictures.
Registration information such as name, password and email address
are required.


\subsubsection{/login}

The direct login area if you are already a registered as an user.
\begin{itemize}
\item Login
\end{itemize}

\subsubsection{/logout}

The direct logout area if you have been previously logged in.
\begin{itemize}
\item Logout
\end{itemize}

\subsubsection{/profile}

A profile might look as followed. At this point the user can make any 
changes to his profile settings, within the range of his permissions. 
This would be eg. loading up an avatar.
\begin{itemize}
\item Setting
\item Avatar
\item Mail
\item Name
\item ...
\end{itemize}

\subsubsection{/user/\{\$id\}}

With the help of this resource you are able to get the profile of
a specific user (which is uniquely identified by its ID). This is
especially useful for viewing all of the galleries of the user.
\begin{itemize}
\item User profile 
\item Overview of public / private galleries
\end{itemize}

\subsubsection{/user/\{\$id\}/galleries/\{\$id\}}

If you want to directly get a specific gallery of an user, you have
to specify the ID of the gallery as a query parameter.
\begin{itemize}
\item Gallery of a single user
\end{itemize}

\subsubsection{/search}

Using the search function, you can get special content. These are selected
with filters.
\begin{itemize}
\item Search
\item Filters
\end{itemize}

\subsubsection{/install}

At the beginning of a session, the installation must be carried out.
\begin{itemize}
\item Installation routine
\end{itemize}

\subsection{Code}

In the following, the program code is explained and illustrated.


\subsubsection{Controllers}

Here is how to recognize the control classes that were implemented and used.

Almost all of these controller classes differ only in their title.
The page title, which is visible in the browser connection is \textquotedbl{}set\textquotedbl{}
at this point. 

For example, the title of 'class-error controller' on \textquotedbl{}404
- Page does not exists!\textquotedbl{} was placed. Another one was described
in the controller.


\subsubsection{Form-Validation}

It was of great importance that measures of validation were implemented in all forms, 
such as the Login or Register forms. 
When logging in it has to be verified, whether the username is available and whether the
associated password is correct.

In order to registrate successfully, an availible username and an unique email address are required. 
Furthermore, it has to be ensured that both passwords entered are the same.

These error handlings are implemented in the class \textquotedbl{}class-user-manager\textquotedbl{}
and shown by the functions ,,validate\_new\_user``, ,,set\_current\_user``
and ,,create\_user through``.

First, a variable is created for each field of the form. If the field is no longer empty after the user's interaction, the value is what the user has entered. Then the entered value is fetched and processed using the POST method. In addition, it is converted by the htmlspecialchars function. This is particularly important in order to maintain an adequate security level. In general this function transforms single characters with the use of selected flags, so that they are not interpreted as false. In this context, it is particularly important for quotation marks, which are converted to prevent misunderstandings as <input value="Mein" Name"> (the value would only be "Mein" and Name" would be considererd as an invalid HTML-Object).


\subsubsection{CSS-Framework Bootstrap}

As a basis and as a simplification, the work was accomplished with the free Bootstrap 3.0 framework. 

This simplified, especially at the beginning, the realization and the rapid progress. 
Bootstrap has the advantage, among others, to be platform independent, but also supports in its 
latest version the \textquotedbl{}responsive design\textquotedbl{}. The site has a dynamic structure, 
regardless of which device was used.

\subsubsection{Functions}
In the following there are declarations about the used functions written in the functions.php.
In addition it should be said that the functionality is only roughly explained. Further descriptions can be found in the source code.

Before the installation process begins, a function checks the compatibility of the installed php-Version. If it is the case, it's important to warn the user that he has an older php-Version or that the MySQLi extension is not started. 

Another function checks, whether or not the installation was already completed before. If this isn't the case it redirects to the installation page. 

To guarantee the safety of the passwords a hash function was implemented. 
For this purpose, the existing \hyperlink{http://www.openwall.com/phpass/}{hash function framework}  was used. 

\subsubsection{Javascript}
It follows a short description of the javascript functions which were implemented. 

First found in the javascript part of the scripts that were taken completely from the specified sources to implement the Fancybox and the slider of the home page for example. These were the contents of complete packages which were available for the free use.

The most important file is the javascript image-uploader.js.
Here, the entire upload of the images and thus the actual functionality is implemented. Furthermore, both Ajax and JSON elements elements are introduced, as was requested in the job position.

First, the images are (this is new in HTML5) selected via a multiple-file upload input and loaded locally using the File Reader to then create a thumbnail. After the upload an AJAX request, which runs on an XML HTTP request level 2 begins.

Password strength function
As an extra to the user a feature was implemented, which determines the strength of the entered password.
\end{document}

